% =====================================================================
% DRAFT — Methods: AIF quality control, features and SQ1 statistics
% Paste the subsections into sections/02_methods.tex. References: methods_aif_refs.bib.
% Numbers come from results/combined_v1 (see aif_qc_followup.md); re-check after re-running the scripts.
% =====================================================================

\subsection{AIF quality control}
\label{sec:aif_qc}

Not every automatically extracted AIF is a valid arterial input: the mask can miss the artery, average
overlapping vessels, or the curve can be distorted by the injection or the acquisition. Because every
perfusion parameter is computed relative to the AIF, an invalid AIF propagates directly into the maps.
Quality control (QC) was therefore applied to every curve (136 unique runs $\times$ 3 mask variants
$=$ 408 curves) before any comparison, in three steps.

\subsubsection{Curve-shape QC}
Each AIF was taken at its native frame times (not the resampled curve, which hides frame-to-frame noise),
baseline-corrected with the median of the pre-arrival frames and normalised to its peak, so that every
threshold is a fraction of the bolus amplitude and independent of dose and display scaling. Six failure rules
were applied in order of how fundamental the defect is: no bolus above the noise, no stable baseline,
truncation (the curve does not wash out), more than one first-pass peak or notch, a late re-rise after
washout, and a noisy tail. A curve that passes all rules is labelled \emph{usable}.

\subsubsection{Failure taxonomy}
To separate segmentation errors from other causes, every failed curve was assigned one class
(Table~\ref{tab:taxonomy}): \emph{wrong structure} (mask scattered over the head, $\geq$3 parts with a spatial
spread $>$200\,px, or touching the image border), \emph{too small} ($<$500\,px), \emph{split segments}
(largest connected part $<$80\,\% of an ICA mask, typically supraclinoid plus cervical ICA),
\emph{projection overlap} (most mask pixels still opaque in the venous phase, or a defect present in one
plane only) and \emph{not segmentation} (the same defect in both planes and all masks, i.e.\ in the bolus or
the acquisition). Geometric rules take precedence because a broken mask makes the cause analysis of the
curve meaningless. The rules were set on, and every class checked against, overlays of the mask on the
minimum-intensity projection.

\subsubsection{Objective QC criteria and trust flag}
The shape QC cannot detect a curve that looks like a bolus but is averaged over the wrong pixels. Each
criterion was therefore defined as one measurable quantity with one threshold and a stated basis
(Table~\ref{tab:qc_criteria}): number of first-pass peaks ($\leq$1), tail fraction (minimum after the peak
relative to the peak; $\leq$0.15, see Section~\ref{sec:truncation}), saturated mask pixels ($\leq$10\,\%, a natural gap in the cohort between 6\,\% and
82\,\%), mask size ($\geq$500\,px; AIF noise scales with $1/\sqrt{N}$), scattered or border masks, and two
\emph{review} criteria (split ICA mask, majority venous mask pixels). A curve is \emph{rejected} if it fails
the shape QC or a reject criterion, flagged for \emph{review} if only a review criterion fails, and
\emph{trusted} otherwise. Only trusted curves enter the statistical analyses. The sensitivity of every
threshold was assessed by recounting at a stricter and a more lenient value.

% NOTE (2026-10-26): this text describes the 0.15 threshold; the pipeline (aif_qc.py TRUNC_LEVEL,
% aif_taxonomy.py TAIL_MAX) still runs with 0.30 until that change is made and everything is re-run.
\subsubsection{Truncation threshold}
\label{sec:truncation}
An AIF is truncated when the first-pass bolus is not captured completely, i.e.\ it has not washed out by the
end of the run~\cite{borst2015}. Truncation lowers the AIF area that normalises CBV and biases MTT; in MR
perfusion, shortening the acquisition lowered CBV by up to 64\,\% and raised MTT by up to 205\,\%~\cite{copen2014}.
No truncation threshold exists for DSA perfusion with perfDSA~\cite{su2025perfdsa}. The threshold was therefore
derived on this cohort following the approach of de~la~Rosa et al.~\cite{delarosa2023}: in runs with a trusted,
completely washed-out AIF, frames were removed one at a time from the end and the perfusion maps recomputed with
the unchanged deconvolution. A shortened run was considered reliable when the median CBV, CBF and MTT deviated
less than 10\,\% from the full run (the reliability bound of~\cite{delarosa2023}); the optimal scan duration is
the shortest duration from which every longer version is reliable. It was expressed as the AIF coverage (time
after the AIF peak; the strongest single predictor of truncation in~\cite{delarosa2023}), as the time after bolus
arrival (the quantity of Lyman et al., who required $\geq$7\,s in DSA~\cite{lyman2026}), and as the AIF tail
fraction reached at that point. Requirements were set at the 95th percentile over the runs (5th percentile for
the tail fraction), so that 95\,\% of runs are reliable. A second simulation lifted only the AIF tail to a fixed
floor while keeping the complete run, which models contrast lingering in the mask (venous overlap, slow
outflow), the situation of the curves labelled truncated in this cohort. The QC tail-fraction threshold was set at
the highest level at which the median CBV, CBF and MTT bias over 56 simulation runs stayed within 10\,\%:
a tail fraction of 0.15 (median bias CBV $-8$\,\%, MTT $-7$\,\%; at the previously used 0.30: $-36$\,\% and
$-30$\,\%). Because de~la~Rosa et al.\ apply the 10\,\% bound per scan, which at 0.15 holds for 62\,\% of the runs
and at 0.10 for 86\,\%, the main comparisons were repeated with a tail threshold of 0.10 as a sensitivity analysis.
Runs became reliable at a median of 6.7\,s after bolus arrival, consistent with the 7\,s reported for DSA by
Lyman et al.~\cite{lyman2026}; because the run length is shared by the curves of a paired comparison (same run
for full versus top, same injection for plane A versus B), run length was reported per run rather than used to
exclude runs.

\subsection{AIF features}
\label{sec:features}
Features were computed on the 250\,ms resampled AIF $c(t)$ after subtracting the baseline $b$ (median of the
first two samples). \emph{Timing}: arrival time (upward crossing of a fraction $f$ of the peak, linearly
interpolated), time to peak (TTP), rise time (TTP $-$ arrival), full width at half maximum (FWHM) and bolus mean
transit time (centroid $-$ arrival). \emph{Amplitude}: peak height, area under the curve (AUC) and first-pass AUC
(from the 10\,\% arrival until the curve falls below 10\,\% of the peak), which excludes the venous and
recirculation tail that the full AUC includes. \emph{Shape-only} features remove the amplitude, so that shape can
be compared between masks and views independent of the amount of contrast: peak-normalised ($c/\mathrm{peak}$:
FWHM, rise time, AUC/peak, up/down ratio, tail fraction) and AUC-normalised ($c/\mathrm{AUC}$: peak/AUC, temporal
SD, skewness).

The arrival fraction $f$ was chosen objectively. Both planes image the same bolus at the same moment, so the A/B
arrival difference measures arrival-time noise; it was nearly constant for $f$ between 0.05 and 0.5 (SD
0.12--0.13\,s, half a frame interval), so precision did not decide. The level must, however, lie above the
baseline noise: $f = 0.15$ is the lowest candidate above the 3$\sigma$ frame noise of every trusted curve and lies
within the conventional 10--20\,\% range. All thresholds are stored in one configuration file, and every feature
is stored per curve together with its QC labels and a data dictionary.

\subsection{Projection differences between planes}
\label{sec:plane_methods}
To explain why the AIF differs between the frontal (A) and lateral (B) view of the same injection, two
candidate sources were quantified per run. \emph{Contrast medium seen by the detector}: the whole-image contrast
signal (sum over all pixels of the subtracted signal, scaled by the pixel area at the isocentre), which in an
ideal projection is the projected iodine mass and identical in both views; dividing the AIF peak by it removes
global scaling (tube voltage, exposure, display processing) and leaves the vessel-specific part.
\emph{Vessel overlap}: the fraction of mask pixels still opaque in the venous phase and the fraction of vessel
pixels in a 15\,px ring around the mask. Planes were compared with paired Wilcoxon signed-rank tests (Holm
correction) and the AIF ratio B/A was related to the candidate factors with Spearman correlation.

\subsection{Statistical analysis: sampling segment (SQ1)}
\label{sec:sq1_methods}
SQ1 asks whether the AIF depends on which ICA segment it is sampled from. For every run, each full-ICA AIF
(\texttt{ica\_full\_model}, \texttt{ica\_full\_grown}) was paired with the supraclinoid \texttt{ica\_top} AIF of
the same run, only when both curves were trusted, so that a pair never contains a rejected curve. The analysis
was done per plane, because the projection changes which part of the ICA each mask covers.

Paired differences (full $-$ top) were taken as plain differences for timing features and dimensionless shape
features, and as log ratios for amplitude and perfusion features, whose error is proportional to their size.
Normality of the differences was assessed with the Shapiro--Wilk test and Q--Q plots, because the paired
$t$-test assumes normally distributed differences, whereas the Wilcoxon signed-rank test only assumes a
symmetric distribution. The test was chosen by a pre-specified rule: paired $t$-test if Shapiro--Wilk
$p \geq 0.05$, otherwise Wilcoxon. Both $p$-values were kept to check that conclusions do not depend on this
choice. $P$-values were Holm-corrected over the features within one comparison and plane. Agreement was
quantified with Bland--Altman analysis (bias and 95\,\% limits of agreement, bias\,$\pm$\,1.96\,SD;
non-parametric 2.5th--97.5th percentile limits when the differences were not normal; ratio features on a log
scale).

To test whether the full-versus-top difference changes direction between views, the difference of differences
($d_B - d_A$) was tested per injection with the Wilcoxon signed-rank test (Holm-corrected). Three candidate
explanations were related to the full-versus-top difference with Spearman correlation: pixel count (log ratio of
mask sizes), vessel overlap (difference in ring vessel fraction and venous fraction between the masks) and
foreshortening of the supraclinoid segment (size and spatial spread of the \texttt{ica\_top} mask).
